\reading{CR-22}{Structured Credit Risk}{STRIKE FIRST}{7}

\begin{crkeybox}[Learning objectives for this reading]
\begin{enumerate}[label=\textbf{\alph*.},leftmargin=2em]
\item Describe common types of structured products.
\item Describe tranching and the distribution of credit losses in a securitization.
\item Describe a waterfall structure in a securitization.
\item Identify the key participants in the securitization process and describe conflicts of interest that can arise in the process.
\item Calculate and evaluate one or two iterations of interim cashflows in a three-tiered securitization structure.
\item Describe the treatment of excess spread in a securitization structure and estimate the value of the overcollateralization account at the end of each year.
\item Explain the tests on the excess spread that a custodian must go through at the end of each year to determine the cash flow to the overcollateralization account and to the equity noteholders.
\item Describe a simulation approach to calculating credit losses for different tranches in a securitization.
\item Explain how the default probabilities and default correlations affect the credit risk in a securitization.
\item Explain how default sensitivities for tranches are measured.
\item Describe risk factors that impact structured products.
\item Define implied correlation and describe how it can be measured.
\item Identify the motivations for using structured credit products.
\end{enumerate}
\end{crkeybox}

\begin{crtrapbox}[Two layers, three numbering schemes, and a duplicated block]
The Crash layer numbers this reading 28, the Lecture layer numbers it 40, and the
duplicate of the Lecture block numbers it 24. The Lecture layer also uses its own
letters, which run one behind the spine from objective c onward: its ``40.g'' is
the spine's h, its ``40.h'' is i, and so on.

\smallskip
The Lecture layer covers spine objectives a, b, d, and h through m. It omits
\textbf{c, e, f and g} --- the waterfall and all three calculation objectives.
Everything below is renumbered to the spine.
\end{crtrapbox}

% =====================================================================
\lo{CR-22 a}{Describe common types of structured products.}

\begin{enumerate}
\item \term{Covered bonds.} Mortgage loans are aggregated into a \term{cover
pool} by which a bond issue is secured. The cover pool \textbf{stays on the
balance sheet} of the bank rather than being sold off-balance sheet. The bank
promises to pay back the borrowed money and interest \emph{no matter how the
mortgages perform}. Unlike typical securitizations, covered bonds do not create a
separate entity, and investors can hold the bank responsible if there are issues.
Investors have priority over other creditors in case of bank default.
\item \term{Mortgage pass-through securities.} In contrast to covered bonds,
these are \textbf{true off-balance sheet} securitizations. Investors receive cash
flows based entirely on the performance of the pool, less associated fees paid to
the servicer. Less default risk, as they are agency MBS, but prepayment risk still
exists.
\item \term{Collateralized mortgage obligations (CMOs).} MBS divided into tranches
with different predetermined conditions. The most basic structure is the
\term{waterfall} or \term{sequential pay} structure, where Tranche 1 receives all
principal and its portion of interest in each period until it is paid off. Other
tranches receive interest until Tranche 1 is paid off, then they receive
principal.
\item \term{Structured credit products.} Backed by risky debt instruments and
divided into tranches with different credit risks. The equity tranche is the
riskiest and the first to bear losses from defaults; if it is wiped out, the next
junior tranche takes the risk. The senior tranche has the highest credit rating
and the lowest chance of writedowns.
\item \term{Asset-backed securities (ABS).} The most general category: any pool of
cash-generating assets can be bundled and tranched. MBS is a specific type of ABS.
Others include CBO, CDO, CLO, CMO, and even complex securities like CDO-squared.
\end{enumerate}

\begin{crkeybox}[Four dimensions on which structured credit products vary]
Underlying collateral; the size and number of tranches; whether the pool is
passive or actively managed; and whether it is \term{revolving}, that is, whether
loan proceeds are reinvested in new assets during a revolving period.
\end{crkeybox}

\begin{crtrapbox}[The covered bond is the odd one out]
It is the only product on the list where the collateral stays on the issuer's
balance sheet and the issuer remains liable regardless of pool performance.
Everything else on the list transfers the performance risk to the investor.
\end{crtrapbox}

% =====================================================================
\lo{CR-22 b}{Describe tranching and the distribution of credit losses in a securitization.}

The capital structure of a securitization refers to the \term{priority} assigned
to the different tranches.

\begin{center}
\footnotesize
\begin{tabular}{L{2.8cm} L{5.8cm} L{5.9cm}}
\toprule
\thead{Tranche} & \thead{Priority} & \thead{Return} \\
\midrule
\term{Senior} & Top of the capital structure; highest priority to receive
  principal and interest &
  Perceived safest, so receives the \emph{lowest} coupon \\
\term{Mezzanine} & Junior to senior; absorbs losses after the equity tranche &
  Relatively high coupon or spread \\
\term{Equity} & Lowest priority; absorbs \emph{first} losses; smallest part of
  the structure &
  Variable return, that is, the residual \\
\bottomrule
\end{tabular}
\end{center}

\subsection*{Credit enhancement}

\begin{center}
\footnotesize
\begin{tabular}{L{3.1cm} L{11.4cm}}
\toprule
\thead{Type} & \thead{Description} \\
\midrule
\term{External} & Buying insurance, or \term{wraps}, from a third party such as a
  monoline insurer \\
\term{Overcollateralization} & \textbf{Hard} credit enhancement: the asset pool
  value exceeds the total bond value, for example 101 mortgages backing 100
  mortgage bonds. Provides protection \emph{at origination} \\
\term{Excess spread} & \textbf{Soft} credit enhancement: the difference between
  cash collected from the assets and payments made to bondholders. Saved in a
  trust and available to make up \emph{future} shortfalls \\
\bottomrule
\end{tabular}
\end{center}
\figcap{Hard versus soft is the distinction to hold, and it maps onto timing.
Overcollateralization is present from day one; excess spread has to be earned
year by year before it protects anything.}

% =====================================================================
\lo{CR-22 c}{Describe a waterfall structure in a securitization.}

A waterfall structure details the distribution of collateral cash flows to the
different classes of bondholders.

\begin{center}
\begin{tikzpicture}[font=\sffamily\footnotesize,
  box/.style={draw, line width=0.8pt, rounded corners=1pt, text width=5.0cm,
              align=center, minimum height=0.78cm}]
\node[box, draw=FRMNavy, fill=PaleNavy] (coll) at (0,0)
  {\scriptsize\textbf{Collateral cash flows}};
\node[box, draw=FRMGrey, fill=FRMSoft, below=4pt of coll] (fees)
  {\scriptsize less fees and defaults};
\node[box, draw=FRMGreen, fill=PaleGreen, below=4pt of fees] (sen)
  {\scriptsize\textbf{Senior} claims paid first};
\node[box, draw=FRMGold, fill=PaleGold, below=4pt of sen] (mez)
  {\scriptsize\textbf{Mezzanine} claims, if senior paid in full};
\node[box, draw=FRMRed, fill=PaleRed, below=4pt of mez] (eq)
  {\scriptsize\textbf{Equity} receives the residual};
\node[box, draw=FRMNavy, fill=PaleNavy, text width=4.0cm,
  below right=0.55cm and 1.9cm of eq] (oc)
  {\scriptsize\textbf{Overcollateralization}\\\scriptsize\textbf{trust account}};
\foreach \a/\b in {coll/fees, fees/sen, sen/mez, mez/eq}
  \draw[-{Stealth[length=5pt]}, FRMNavy, line width=0.9pt] (\a) -- (\b);
\draw[-{Stealth[length=5pt]}, FRMNavy, line width=0.9pt] (eq.east) -| (oc.north);
\node[font=\scriptsize, align=left, text=FRMNavy, fill=white, inner sep=2pt,
  anchor=west] at ($(eq.east)+(0.25,0.42)$)
  {\scriptsize excess above the OC trigger is diverted};
\draw[-{Stealth[length=6pt]}, FRMGrey, line width=1pt]
  (-3.6,0.45) -- (-3.6,-3.55)
  node[midway, left, font=\scriptsize\itshape, text=FRMGrey, rotate=90,
  anchor=south, yshift=4pt] {order of payment};
\end{tikzpicture}
\end{center}
\figcap{The equity tranche typically receives the residual cash flows once the
senior and mezzanine investor claims are satisfied. If the cash flows to equity
holders exceed the overcollateralization trigger, the excess is diverted to a
trust account. Fees and defaults reduce the net cash flows available for
distribution before any tranche is paid.}

% =====================================================================
\lo{CR-22 d}{Identify the key participants in the securitization process and describe conflicts of interest that can arise in the process.}

\begin{center}
\footnotesize
\begin{tabular}{L{2.9cm} L{5.6cm} L{6.0cm}}
\toprule
\thead{Participant} & \thead{Role} & \thead{Conflict of interest} \\
\midrule
\term{Originator / sponsor} & Funds the original loans and may supply most of the
  collateral & --- \\
\term{Underwriter} & Structures the issue, engineers the tranches, sells the
  bonds. Warehouses the collateral and faces marketing and collateral-value risk &
  --- \\
\term{Credit rating agencies} & Provide credit ratings, influencing the size of
  tranches; may require enhancements &
  Conflicted between generating profit and allocating larger portions of the
  capital structure to lower-interest-paying senior notes \\
\term{Servicer} & Collects and distributes collateral cash flows and fees;
  provides liquidity on late payments and handles defaults &
  May \textbf{delay foreclosure to increase fees}, while investors want quick
  resolutions \\
\term{Manager} \newline (actively managed pools) & Monitors the credit quality of
  the collateral &
  May lack clear incentives to monitor continually \\
\term{Custodians and trustees} & Administrative: verify documents, disburse and
  transfer funds between accounts & --- \\
\bottomrule
\end{tabular}
\end{center}
\figcap{Three of the six participants carry a named conflict, and each conflict
runs in a different direction: the rating agency's distorts the capital structure,
the servicer's distorts the timing of recoveries, and the manager's is a conflict
of \emph{omission}.}

% =====================================================================
\lo{CR-22 e}{Calculate and evaluate one or two iterations of interim cashflows in a three-tiered securitization structure.}

\gapbadge

{\footnotesize\color{FRMGrey}This objective, and objectives f and g, are
\emph{calculate} and \emph{estimate} verbs. The Lecture layer omits all three
entirely. The Crash layer gives the mechanics in prose and then instructs:
``Please Solve an example from the book for this.'' The structure, assumptions and
worked iterations below are written from the GARP curriculum.}

\begin{crgapbox}[What the objectives ask for: the structure and its assumptions]
\textbf{The structure.} Underlying debt instruments of \$100 million with a coupon
of Libor $+$ 350 bps, funded by an equity note of \$5 million, mezzanine debt of
\$10 million at Libor $+$ 500 bps, and senior debt of \$85 million at Libor $+$
50 bps.

\smallskip
\textbf{The assumptions.} The swap curve is flat at 5\%. All loans and liabilities
mature in five years, and all coupons and loan interest payments are annual and
occur at year end. A defaulted loan pays \emph{no} interest for the year in which
it defaults or after; there is no partial interest. Recovery is 40\%, and the
recovery amount is paid into the overcollateralization account rather than through
the waterfall --- otherwise the equity would perversely benefit from defaults. The
overcollateralization account earns the money market rate of 5\%.

\smallskip
\textbf{The overcollateralization mechanism.} Instead of letting all the excess
spread flow to the equity note, up to \$1,750,000 per year is diverted to the
overcollateralization account. Funds in the account are used to pay bond interest
if collateral interest is insufficient. Any remaining funds are released to the
equity tranche \emph{only at maturity}.
\end{crgapbox}

\begin{crexbox}[The structure with no defaults]
Compute the annual cash flows and the excess spread assuming no defaults, and
confirm the attachment points.
\tcblower
\begin{align*}
\text{Collateral} &= (0.050 + 0.0350)\times \$100{,}000{,}000 = \$8{,}500{,}000 \\
\text{Mezzanine} &= (0.050 + 0.0500)\times \$10{,}000{,}000 = \$1{,}000{,}000 \\
\text{Senior} &= (0.050 + 0.0050)\times \$85{,}000{,}000 = \$4{,}675{,}000
\end{align*}
The bond coupon due, $B$, is $\$1{,}000{,}000 + \$4{,}675{,}000 =
\$5{,}675{,}000$, a constant for every year.

\smallskip
\textbf{Excess spread with no defaults:}
\[ \$8{,}500{,}000 - \$5{,}675{,}000 = \$2{,}825{,}000 \]

\textbf{Attachment points.} The mezzanine attachment point is the equity tranche
as a share of the pool, $5/100 = \mathbf{5\%}$. The senior attachment point is
equity plus mezzanine, $15/100 = \mathbf{15\%}$. There is initially no
overcollateralization.

\smallskip
\textbf{Weighted average spread on the debt tranches:}
\[ \frac{(10 \times 500) + (85 \times 50)}{95} = \frac{9{,}250}{95}
   = \mathbf{97.4\ \text{bps}} \]
The junior bond has a much wider spread than the senior and much less credit
enhancement.
\end{crexbox}

\begin{crexbox}[Two iterations of interim cash flows at a 5\% default rate]
The pool holds 100 loans of \$1 million each. The default rate is 5\% annually,
defaulted collateral is not replaced, and defaults are assumed to occur evenly.
Track years 1 and 2.
\tcblower
\textbf{Year 1.} Defaults $= 5\% \times 100 = 5.00$ loans, leaving 95.00
surviving.
\begin{align*}
\text{Loan interest } B_1 &= 95.00 \times \$1{,}000{,}000 \times 0.085
  = \$8{,}075{,}000 \\
\text{Excess spread} &= \$8{,}075{,}000 - \$5{,}675{,}000 = \$2{,}400{,}000
\end{align*}
Excess spread exceeds the \$1,750,000 cap, so \$1,750,000 is diverted to the
overcollateralization account and $\$2{,}400{,}000 - \$1{,}750{,}000 =
\$650{,}000$ is paid to equity.

Recovery $= 5.00 \times \$1{,}000{,}000 \times 40\% = \$2{,}000{,}000$, also paid
into the account.
\[ \text{OC balance, end of year 1} = \$1{,}750{,}000 + \$2{,}000{,}000
   = \mathbf{\$3{,}750{,}000} \]

\smallskip
\textbf{Year 2.} Defaults $= 5\% \times 95.00 = 4.75$ loans, leaving 90.25
surviving. \emph{The annual number of defaults falls because the pool has
shrunk.}
\begin{align*}
\text{Loan interest } B_2 &= 90.25 \times \$1{,}000{,}000 \times 0.085
  = \$7{,}671{,}250 \\
\text{Excess spread} &= \$7{,}671{,}250 - \$5{,}675{,}000 = \$1{,}996{,}250
\end{align*}
Still above the cap, so \$1,750,000 is diverted and \$246,250 goes to equity.

Recovery $= 4.75 \times \$1{,}000{,}000 \times 40\% = \$1{,}900{,}000$.

The existing balance earns the money market rate:
\[ \text{OC balance, end of year 2}
   = \$3{,}750{,}000 (1.05) + \$1{,}750{,}000 + \$1{,}900{,}000
   = \mathbf{\$7{,}587{,}500} \]
\end{crexbox}

% =====================================================================
\lo{CR-22 f}{Describe the treatment of excess spread in a securitization structure and estimate the value of the overcollateralization account at the end of each year.}

The three-tiered waterfall has scheduled payments to the senior and mezzanine
tranches. The equity tranche receives cash flows only if excess spread is positive,
that is, if interest collected exceeds interest owed to senior plus mezzanine.

\begin{crfmlbox}[Excess spread and the overcollateralization account]
\[ \text{excess spread}_t \;=\; B_t - B \]
\[ \mathrm{OC}_t \;=\; \min\bigl(\max(\text{excess spread}_t,\,0),\; K\bigr) \]
\[ \text{OC balance}_t \;=\; \text{OC balance}_{t-1}(1+r)
   \;+\; \mathrm{OC}_t \;+\; \text{recovery}_t \]
\vk{$B_t$ = aggregate loan interest received by the trust at the end of year $t$;
$B$ = bond coupon interest due to both the junior and senior bonds, a constant;
$K$ = the maximum amount divertible annually, here \$1,750,000; $r$ = the money
market rate. The account increases from recovery of defaulted assets and from the
diversion of spread, and earns the money market rate on its balance.}
\end{crfmlbox}

If excess spread is \emph{negative} --- interest collected is less than interest
owed to senior plus mezzanine --- the OC account will use all of its available
funds until depleted.

\begin{crkeybox}[Where the structure breaks]
With 100 loans at \$1 million paying 8.5\%, excess spread turns negative once
fewer than
\[ \frac{\$5{,}675{,}000}{\$1{,}000{,}000 \times 0.085} = 66.76 \]
loans survive. Below that point the overcollateralization account starts being
drawn down rather than built up, and the equity has already been receiving nothing
for some time.
\end{crkeybox}

% =====================================================================
\lo{CR-22 g}{Explain the tests on the excess spread that a custodian must go through at the end of each year to determine the cash flow to the overcollateralization account and to the equity noteholders.}

\begin{center}
\begin{tikzpicture}[font=\sffamily\footnotesize,
  t/.style={draw=FRMNavy, fill=PaleNavy, line width=0.8pt, rounded corners=1pt,
            text width=5.4cm, align=center, minimum height=0.8cm},
  o/.style={draw, line width=0.8pt, rounded corners=1pt, text width=4.6cm,
            align=center, minimum height=1.25cm}]
\node[t] (q1) at (0,0) {\scriptsize Is excess spread $> 0$?};
\node[o, draw=FRMRed, fill=PaleRed] (no1) at (7.4,0)
  {\scriptsize\textbf{No}\\[1pt]\scriptsize OC account pays the shortfall on the bonds, until depleted. Equity receives nothing};
\node[t, below=1.35cm of q1] (q2) {\scriptsize Does it exceed the cap $K$?};
\node[o, draw=FRMGold, fill=PaleGold] (no2) at (7.4,-1.85)
  {\scriptsize\textbf{No}\\[1pt]\scriptsize all of it is diverted to the OC account. Equity receives nothing};
\node[o, draw=FRMGreen, fill=PaleGreen] (yes2) at (7.4,-3.85)
  {\scriptsize\textbf{Yes}\\[1pt]\scriptsize $K$ is diverted to the OC account; the excess over $K$ is paid to equity};
\draw[-{Stealth[length=5pt]}, FRMRed, line width=0.9pt] (q1) -- (no1);
\draw[-{Stealth[length=5pt]}, FRMNavy, line width=0.9pt] (q1) -- (q2)
  node[midway, right, font=\scriptsize, text=FRMNavy] {yes};
\draw[-{Stealth[length=5pt]}, FRMGold, line width=0.9pt] (q2) -- (no2);
\draw[-{Stealth[length=5pt]}, FRMGreen, line width=0.9pt] (q2.south) |- (yes2.west);
\end{tikzpicture}
\end{center}
\figcap{The two tests the custodian runs at each year end, and the three outcomes
they produce. Note that equity receives cash in exactly \emph{one} of the three
branches, and that any balance left in the account is released to equity only at
maturity --- so the equity holder is last in line twice over.}

% =====================================================================
\lo{CR-22 h}{Describe a simulation approach to calculating credit losses for different tranches in a securitization.}

Defaults happen at different times, loans have different default risks, and loans
can be correlated, so a simulation approach is needed to address this complexity.

\begin{enumerate}
\item \term{Estimate critical parameters.} Determine default intensity from market
spread data, and estimate correlations between loans. Correlation is challenging
to estimate precisely because of a lack of usable market data, so \term{copulas}
are used.
\item \term{Generate default time simulations.} Use simulations to create various
scenarios as to \emph{when} and \emph{if} each security defaults.
\item \term{Compute portfolio credit losses.} Calculate credit losses for each
tranche across all simulation scenarios to determine the frequency and timing of
credit losses.
\end{enumerate}

The credit losses can be lined up to assess the impact on the capital structure.
The \emph{tail} of the distribution identifies the \term{credit VaR} for each
tranche.

% =====================================================================
\lo{CR-22 i}{Explain how the default probabilities and default correlations affect the credit risk in a securitization.}

\subsection*{Increasing default probability, holding correlation constant}

\begin{center}
\footnotesize
\begin{tabular}{L{3.6cm} c c}
\toprule
\thead{Tranche} & \thead{Mean value} & \thead{Credit VaR} \\
\midrule
Equity tranche & $\downarrow$ & $\downarrow$ \\
Mezzanine tranche & $\downarrow$ & $\uparrow$ then $\downarrow$ \\
Senior tranche & $\downarrow$ & $\uparrow$ \\
\bottomrule
\end{tabular}
\end{center}
\figcap{Value falls everywhere --- that column is uniform and uninteresting. The
credit VaR column is where the discrimination is: it falls for equity (less
variation in returns, since writedowns become near-certain), rises for senior
(more variation), and for mezzanine rises at low default levels like a senior bond
and then falls at high default levels like equity.}

\subsection*{Increasing correlation}

\begin{center}
\footnotesize
\begin{tabular}{L{3.2cm} L{11.2cm}}
\toprule
\thead{Tranche} & \thead{Effect of \emph{increasing} correlation on value} \\
\midrule
\term{Senior} & \textbf{Falls.} Higher correlation makes extreme pool-wide losses
  possible, which is the only thing that can reach a senior tranche. Equivalently,
  \emph{decreasing} correlation increases senior value \\
\term{Equity} & \textbf{Rises.} The possibility of zero or very few credit losses
  increases with high correlation \\
\term{Mezzanine} & \textbf{Mixed.} At \emph{low} default rates increasing
  correlation decreases mezzanine value; at \emph{high} default rates it increases
  it \\
\bottomrule
\end{tabular}
\end{center}
\figcap{Correlation is a bet on dispersion, and the two ends of the capital
structure take opposite sides of it. At low correlation defaults are predictable
and close to their expected value, which assures the equity tranche of writedowns
and protects the senior. At high correlation outcomes become extreme in both
directions, which is good for equity and bad for senior. Mezzanine sits in between
and flips.}

% =====================================================================
\lo{CR-22 j}{Explain how default sensitivities for tranches are measured.}

For all increases in default rates, equity values fall and bond losses rise. That
is universally true. When correlations change from higher or lower levels, we see
a \term{convex, nonlinear} impact, which affects each tranche differently and
gives each a different risk profile. A granular measure of sensitivity is
therefore needed.

\begin{crdefbox}[Default '01]
The \term{default '01} measures the impact of an increase of one basis point in
the default probability. To compute it, increase and decrease the default
probability by 10 bps and revalue each tranche at these new values.
\end{crdefbox}

\begin{crkeybox}[The gamma of the default '01]
The \emph{gamma} of the default '01 --- the rate of change in the rate of default
sensitivity --- is highest, that is most sensitive, in \term{low-correlation}
environments. As losses become large, the '01 has less of an effect.

\smallskip
This is the same idea as a deeply in-the-money option going even deeper in the
money: it has an almost one-for-one impact on the price of the asset, and the
gamma is low.
\end{crkeybox}

% =====================================================================
\lo{CR-22 k}{Describe risk factors that impact structured products.}

\begin{enumerate}
\item \term{Systematic risk.} Even a well-diversified pool of loans cannot
eliminate systematic risk. High correlations can lead to severe damage to the
portfolio, especially for \emph{senior} tranches, increasing the likelihood of
writedowns.
\item \term{Tranche thinness.} Equity and mezzanine tranches are relatively thin,
so the 95\% and 99\% credit VaR sit close together. A small change in value can
result in a large loss, and breached tranches are more likely to incur significant
losses.
\item \term{Loan granularity.} A pool with fewer but larger loans, for example
commercial mortgages, is more susceptible to unexpected defaults than a pool with
many smaller loans, for example residential mortgages. Reducing the number of
loans while increasing their size raises the probability of tail events.
\end{enumerate}

% =====================================================================
\lo{CR-22 l}{Define implied correlation and describe how it can be measured.}

\begin{crdefbox}[Implied correlation]
Implied correlation is a very similar concept to the implied volatility of an
equity option. By using observed market prices and a pricing model for tranches,
the unique implied correlation can be derived: it is the correlation level that
must be plugged into the model to reproduce the market price of a tranche.
\end{crdefbox}

The process involves three steps:

\begin{enumerate}[label=\alph*)]
\item Use the observable CDS term structure to extract risk-neutral default
probabilities and recovery rates.
\item Assuming constant pairwise correlation, feed the market prices for the
respective tranches together with the default estimates into a copula.
\item Calculate the risk-neutral implied correlation, that is, the base
correlation, per tranche.
\end{enumerate}

\begin{crkeybox}[Correlation skew]
Implied correlation estimates may vary among tranches, leading to a
\term{correlation skew}. The distinction between compound (tranche) correlation
and base correlation, and the smile-versus-skew shapes they produce, is at
\textbf{CR-13 g}.
\end{crkeybox}

% =====================================================================
\lo{CR-22 m}{Identify the motivations for using structured credit products.}

\begin{center}
\footnotesize
\begin{tabular}{L{3.2cm} L{11.2cm}}
\toprule
\thead{Party} & \thead{Motivation} \\
\midrule
\term{Loan originators} & Lower cost funding by selling loans into a
  securitization pool, especially if the loans are high quality and
  well-diversified \\
\term{Investors} & Access to a diversified pool of loans they would not otherwise
  reach, and the ability to \emph{choose} a desired level of risk and return by
  selecting different tranches \\
\bottomrule
\end{tabular}
\end{center}
